\documentclass[10pt,a4paper]{amsart}
\usepackage[margin=19mm]{geometry}
\usepackage{amsmath,amssymb,mathtools}
\usepackage[scr=rsfs,cal=euler]{mathalfa}
\usepackage{xfrac}
\usepackage[unicode,hidelinks,bookmarks=false]{hyperref}
\newcommand{\ie}{i.e.\ }
\newcommand{\cf}{cf.\ }
\newcommand{\resp}{respectively\ }
\newcommand{\Subsection}{\subsection}
\providecommand{\nobreakdash}{\nobreak\hbox{-}}
\setlength{\emergencystretch}{2em}
\multlinegap0pt
\usepackage{amssymb}
\usepackage{tikz}
\usepackage{enumerate}
\usepackage{dsfont}
\newtheorem{lem}{Lemma}
\newtheorem{prop}{Proposition}
\usepackage{cleveref}
\crefname{lem}{Lemma}{Lemmas}
\crefname{prop}{Proposition}{Propositions}
\newcommand{\F}{\mathbb F}
\newcommand{\HH}{\mathcal{H}}
\title{On Reed-Muller subcodes, Grassmannian partitions and sum-free functions}
\author{Philipp Heering, Christian Kaspers, Vladislav Taranchuk}
\date{Source: arXiv:2605.22958v1 (CC BY 4.0)}
\begin{document}
\maketitle

\noindent\emph{Source and licence.} On Reed-Muller subcodes, Grassmannian partitions and sum-free functions. Version arXiv:2605.22958v1 (CC BY 4.0), distributed by arXiv under the Creative Commons Attribution 4.0 International licence, http://creativecommons.org/licenses/by/4.0/, as recorded in the arXiv licence metadata for this exact version and shown on the abstract page https://arxiv.org/abs/2605.22958v1. Full licence notice: \texttt{licenses/arxiv-2605.22958-CC-BY-4.0.txt}.\par\smallskip

\noindent\textit{Editorial scope.} Counted unit: the article's prop 5.1 (source0.tex lines 718--720). The article's complete original proof of it is delivered with it (source0.tex lines 722--766, one \texttt{proof} environment of 45 lines). No mathematics is written, repaired or re-derived: the statement and the proof are the article's own lines, in the article's own notation, and no retained line is substituted or re-flowed.  Retained uncounted as the source-local context the proof needs: lines 381--383; lines 671--677.  Nothing was omitted from any retained span, so the package carries no omission record.  No retained line contains a citation, so the wrapper carries no bibliography and no bibliography entry.  No retained line contains a figure environment, an image-inclusion command or drawing code, so no image is delivered and the package declares no asset dependency.  Editorial formatting only, in addition: an AMS \texttt{amsart} preamble that restates the article's own declarations and macros used by the retained lines, namely the environments \texttt{lem}, \texttt{prop} on separate counters, together with the standard packages those lines need.  The retained environments print in this document as Lemma 1, Proposition 1 in order of appearance, whereas the article prints the same items with the numbering of its own full document.  The complete native source of the article as distributed in the arXiv e-print for this exact version is bundled unchanged as \texttt{sources/201172/source0.tex}.
\begin{lem}\label{L: kth-order coloring}
    Let $k < n$ be positive integers and $F$ be a $k$th-order sum-free $(n, m)$-function. Then for any two distinct $k$-flats $A_1, A_2 \subset \F_2^n$ that intersect in a $(k-1)$-flat, we have $\omega_F(A_1)\neq \omega_F(A_2)$. 
\end{lem}

\begin{lem} \label{L: special condition} 
    Let $F$ be an $(n,m)$-function, let $k\geq 2$ and let $x, y \in \F_{2}^n$. Let $U$ be a $(k-1)$-space of $\F_2^n$. If $F$ has algebraic degree $k$, then $$\omega(U)+\omega(x+U)+\omega(y+U)=\omega(x+y+U).$$
    Furthermore, if $F$ is $(k-1)$th-order sum-free, then 
    $$
    \omega(U)+\omega(x+U)+\omega(y+U) \neq 0.
    $$
\end{lem}

\begin{prop} \label{P: constructive Prove of Them 5.1} 
   Let $F$ be an $(n,m)$-function that is $\{k-1,k\}$-order sum-free and of algebraic degree $k$ for $2\leq k<n$. Then $c_F$ induces a valid coloring on $J_2(n+1,k)$.
\end{prop} 

\begin{proof}
Let $U$ be a $k$-space of $\F_2^{n+1}$. As $F$ is $\{k-1,k\}$-order sum-free, we always have  $c_F(U)\ne0$.
   It remains to be proven that $k$-spaces that intersect in dimension $k-1$ have different colors, i.e. $c_F$ assigns them different values.
There are four possible cases that have to be checked; we treat them all separately. In all cases, let $U_1,U_2$ be $k$-spaces of $\F_2^{n+1}$ with $\dim(U_1\cap U_2)=k-1$.

    \noindent \underline{Case 1:} Assume that $U_1,U_2\subseteq \HH$. Then 
    $$
    c_F(U_1) + c_F(U_2) = \Omega(U_1) + \Omega(U_2) =\omega(U_1) + \omega(U_2).
    $$
 \cref{L: kth-order coloring} implies that the quantity above is non-zero.
  

    \noindent \underline{Case 2:}  Assume that $U_1\subseteq \HH$  and  $U_2\not \subseteq \HH$. Then $U_1 \cap U_2 = W \subseteq \mathcal{H}$. 
    Then, there exist  $a,b\in \F_2^n$, such that 
    \begin{align*}
         U_1 &= W \cup (W + (a, 0)), \\
         U_2 &= W \cup (W+(b, 1)).
     \end{align*}
    By definition, we have 
    \begin{align*}
         c_F(U_1) + c_F(U_2)  &=\Omega(U_1) + \Omega(U_2\setminus \HH) \\
      &= \omega(W) + \Omega(W +(a, 0)) + \Omega(W +(b,1))\\
      &= \omega(W) + \omega(W +a) + \omega(W +b).
     \end{align*}
By Lemma \ref{L: special condition}, the quantity above is non-zero, and hence $c_F(U_1) \neq c_F(U_2)$.
    

     \noindent \underline{Case 3:} Assume that $U_1\cap U_2\subseteq \HH$ and $U_1,U_2\not\subseteq \HH$. Then we have
     $$
     c_F(U_1) +c_F(U_2) = \Omega(U_1\setminus \HH) + \Omega(U_2 \setminus \HH).
     $$
     Since $\dim(U_1 \cap U_2) = k-1$, then $(U_1 \setminus \HH) \cup (U_2\setminus \HH)$ is a $k$-flat, the projection on $\HH$ is also a $k$-flat and hence the sum above cannot be zero, as $F$ is $k$th-order sum-free. 

     \noindent \underline{Case 4:} Assume that $U_1,U_2\not\subseteq \HH$ and $U_1\cap U_2\not\subseteq \HH$. Let $W:= U_1 \cap U_2 \cap \HH$. Then, there exist  $a, b, c \in \F_2^n$ with distinct cosets modulo $W$, such that
     \begin{align*}
         U_1 &= W \cup (W + (a, 1)) \cup (W+ (b, 1)) \cup (W + (a+b, 0)), \\
         U_2 &= W \cup (W+(a, 1)) \cup (W + (c, 1)) \cup (W+ (a+c, 0)).
     \end{align*}
     Hence, 
     \begin{align*}
     c_F(U_1) + c_F(U_2)& = \Omega(W + (a, 1)) + \Omega(W + (b, 1)) + \Omega(W + (a, 1)) + \Omega(W + (c, 1)) \\
     &= \Omega(W + (b, 1)) + \Omega(W + (c, 1)).
     \end{align*}
     Since $(W + b)\cup (W + c)$ is a $(k-1)$-flat of $\HH$ and $F$ is $(k-1)$th-order sum-free, this quantity is non-zero.
 \end{proof}

\end{document}
